\documentclass[11pt,a4paper]{article}
\usepackage[T1]{fontenc}
\usepackage[margin=27mm]{geometry}
\usepackage[english]{babel}
\usepackage{lmodern,microtype,amsmath,booktabs,tabularx}
\usepackage{csquotes}
\usepackage[style=apa,backend=biber,natbib=true]{biblatex}
\addbibresource{references.bib}
\usepackage[hidelinks]{hyperref}
\title{Alexithymia and Emotional Expression in AI-Assisted Music Creation: Expressive Scaffolding and Emotional Clarity}
\author{Mateusz Pomianek}
\date{}
\begin{document}
\maketitle
\begin{abstract}
Generative music tools can offer emotionally meaningful language and musical choices to users who find it difficult to articulate feelings. Such expressive assistance may improve the perceived fit of an artefact without improving independent emotional understanding. This prospective protocol specifies a randomised, parallel-group online experiment in which Polish-speaking adults assemble a short musical sketch from identical AI-generated instrumental options, with or without access to pre-generated emotion-language suggestions. The primary estimand is modification of the assignment effect on a composite expressive-success score (perceived fit when a meaningful sketch exists; zero when no meaningful sketch exists) by baseline difficulty describing negative feelings, measured with the Polish Perth Alexithymia Questionnaire. Secondary outcomes distinguish immediate unaided emotional clarity, perceived agency and distress. The design isolates availability of a bounded semantic scaffold rather than the effect of generative music or clinical treatment. A prespecified regression, missing-outcome sensitivity analysis, manipulation checks and sample-size sensitivity calculation are provided. No human-participant data have been collected or analysed.
\end{abstract}
\noindent\textbf{Keywords:} alexithymia; emotional expression; emotional clarity; affect labelling; generative AI; music psychology

\section{Introduction}
Alexithymia is a multidimensional trait involving difficulty identifying and describing feelings and a tendency towards externally oriented thinking. The attention--appraisal account distinguishes attention to emotional states from appraisal of their meaning \parencite{preece2018}. The Polish Perth Alexithymia Questionnaire (PAQ) permits separate assessment of positive and negative feelings and has initial evidence of factorial validity and reliability in Polish adults \parencite{larionow2022}. Self-reported alexithymia should not be equated with objective emotional awareness, clinical diagnosis or absence of emotional experience \parencite{lane2021}.

Emotion regulation encompasses processes serving a person's emotional goals rather than a single direction of affective change \parencite{gross2015}. Music may serve distinct regulatory functions \parencite{saarikallio2007}. Human--AI songwriting introduces additional operations of selection, revision and delegation \parencite{huang2020}. Emotion-language assistance could reduce the demands of generating a description, especially for people reporting difficulty describing feelings. Yet a fluent generated description may reflect the system's competence rather than the user's independent understanding. Evidence on affect labelling \parencite{lieberman2007} motivates this mechanism but cannot establish effects of assisted music creation, nor does the present behavioural design test neural mechanisms.

\section{Theoretical contribution and hypotheses}
The proposed extension of AIM-ER separates three constructs: (i) trait difficulty identifying or describing feelings, (ii) perceived congruence between one's intended expression and a constructed musical sketch, and (iii) unaided understanding of one's present feelings after the scaffold is removed. Expressive fit and clarity can dissociate. A compensatory account predicts stronger benefits of language suggestions among people with higher description difficulty. A dependence account instead predicts limited transfer to unaided understanding or even poorer fit when suggestions are incongruent. These competing accounts are falsifiable and do not presuppose benefit.

\begin{description}
\item[H1 (primary)] The assignment effect of semantic scaffolding on composite expressive success increases with baseline negative difficulty describing feelings (N-DDF). A two-sided interaction test will permit evidence of the opposite direction.
\item[H2 (secondary)] Greater baseline negative difficulty identifying feelings (N-DIF) is associated with lower post-task unaided emotional clarity, conditional on baseline clarity and assignment. This is a trait association, not a causal effect.
\item[H3 (secondary)] Assignment may improve immediate unaided clarity after the scaffold is removed. The average effect will be estimated separately from expressive fit; a two-sided test allows either direction.
\end{description}
No inference of improved emotional ability will be based solely on H1. A precise null or reversed H1 interaction would weaken the compensatory account for this implementation. A positive H1 with imprecise H3 is not evidence of absence of transfer. Agency, perceived dependence, distress and delayed clarity are exploratory.

\section{Methods}
\subsection{Design and recruitment}
This independent experiment complements, rather than duplicates, the episode-anchored observational analysis of motivations in P03/S02. Adults aged at least 18, able to read Polish and hear the stimuli on their device, will be recruited online using neutral advertising. Musical expertise is not required; prior generative music experience and recruitment channel will be recorded. No clinical enrichment is planned. Participants select a mild, non-traumatic everyday negative experience or a fictional analogue; no narrative disclosure is required. The selected mode and baseline affect are recorded before randomisation. The target population is consenting eligible adults completing baseline measures, not all users of music generators.

After baseline, server-side concealed 1:1 allocation assigns access to a bank of emotion-language suggestions or a matched interface without suggestions. Both groups receive the same instructions, time budget, instrumental options, editing functions and opportunity to select or reject an output. Suggestions are fixed, screened AI-generated phrases, with an explicit option to reject all of them. They do not contain personal material. All participants can independently write a short title or lyric fragment locally and combine it with a selected AI-generated instrumental excerpt; the task is an assembly-and-selection proxy, not unrestricted live text-to-audio generation. The musical materials and interface will be versioned and frozen after piloting. Listening duration, available options and interface latency are standardised. Participants cannot be blinded to availability of suggestions; neutral framing and identical assessment procedures mitigate demand effects. Analysts use masked group labels until the analysis script is frozen.

\subsection{Procedure}
Baseline (T0): consent, eligibility, PAQ, device check, musical experience, initial distress, affect, clarity and unaided expression ease. Participants choose the everyday or fictional scenario before assignment. During the task, they assemble one sketch within a fixed eight-minute window. Completion is not an eligibility criterion after randomisation. T1: rate expressive fit of the attempted sketch, including unsuccessful attempts; then close the creation interface, independently rate present clarity and distress, and report agency and task demand. A participant who explicitly confirms that no meaningful sketch exists has a structural zero on the composite, not an observed fit rating. If completion status or an applicable fit rating is unknown, the composite is missing. A technical interruption or withdrawal alone never implies a structural zero. T2 (optional, 24 hours): unaided clarity and perceived adverse consequences. No further intervention is assigned. Logs retain only coarse operation counts and completion flags, not prompts, lyrics or audio uploads.

\subsection{Operationalisation and measurement gates}
The full Polish PAQ is preferred to its short form because facet-level assessment is essential. N-DDF is the primary moderator and N-DIF the secondary predictor; use the authorised item mapping and sum scoring, without dichotomising participants or inventing a clinical threshold. The primary analysis requires complete baseline N-DDF items before randomisation; N-DIF completeness is required only for H2. Report score distributions, floor/ceiling frequencies and omega with bootstrap confidence intervals where item-level modelling is appropriate. The measure will not be modified without a separate validation plan. In the Polish validation sample (\(N=1008\)), omega for N-DDF was .90 (95\% CI [.89, .91]) and for N-DIF .89 ([.88, .90]) \parencite{larionow2022}. These estimates describe the validation sample, not the reliability of the present study. The short PAQ cannot substitute for the full instrument when facet scores are needed.

Expressive fit is a task-specific 0--10 rating of how well a meaningful attempted sketch represents the intended emotion, anchored from no fit to complete fit. The primary outcome is a composite of production and fit: verified absence of a meaningful sketch is assigned zero, while an applicable rating retains its observed value. This operational definition estimates overall expressive success, not fit conditional on completion. Completion proportions and applicable fit distributions are reported separately. It is not a validated clinical scale. Baseline unaided expression ease is a separate 0--10 task-specific covariate. Immediate emotional clarity is a separate anchored 0--10 self-report of understanding one's present feelings, assessed before assignment and again without access to suggestions; the same wording and anchors are used at both times. The S-DERS clarity subscale is a candidate convergent measure, contingent on verified Polish adaptation, licensing and adequate reliability \parencite{lavender2017}. The primary endpoint must pass cognitive interviews and feasibility checks before registration; if comprehensibility or range is inadequate, revise and register before confirmatory recruitment. No claim of performance-based emotional awareness is made. A blinded, independently scored emotional-description task is a possible future study, subject to privacy review and validated scoring.

\subsection{Sample size and feasibility}
The primary interaction is expressed in residual-standard-deviation units of composite expressive success per one baseline moderator SD, comparing assigned arms. Under balanced allocation, centred treatment coding and an independent standardised moderator, its planning standard error is approximately \(2/\sqrt{n}\) in residual-SD units. A two-sided normal approximation at alpha .05 gives about 80 percent power for an interaction of .20 at \(n=800\) observed outcomes and .26 at \(n=480\). These are sensitivity scenarios, not literature-derived effects or empirical estimates. The proposed cap is 1,000 randomised participants, anticipating up to 20 percent unavailable primary outcomes. A separate feasibility pilot of approximately 30--50 adults will inform outcome range, recruitment and availability, not estimate the confirmatory interaction for sample-size selection. Before registration, specify a smallest meaningful interaction in raw outcome units, check feasible recruitment, and evaluate the supplied Monte Carlo scenarios for bounded scores, unequal residual variance, structural zeros and attrition. The simulation parameters are hypothetical; they are not pilot estimates. If 800 observed outcomes are implausible, freeze a smaller target with explicitly reduced sensitivity or designate the work exploratory. Stop at the fixed randomisation cap or preregistered recruitment date, not statistical significance.

\section{Statistical analysis}
\subsection{Estimand and primary model}
Let \(A_i\in\{-.5,.5\}\) denote assignment, \(D_i\) the baseline N-DDF score centred and divided by its SD among all randomised participants, \(E_i\) baseline expression ease and \(C_i\) the prespecified baseline covariates (initial clarity, distress, prior music-generation experience and scenario mode). Fit
\[
Y_i=\beta_0+\beta_1 A_i+\beta_2 D_i+
\beta_3 A_iD_i+\beta_4 E_i+\gamma^\top C_i+\epsilon_i.
\]
\(Y_i\) is composite expressive success at T1, and \(\beta_3\) is the linear effect modification in raw 0--10 points per moderator SD. Use OLS with HC3 standard errors, two-sided alpha .05 and a 95 percent confidence interval. With incomplete outcomes the prespecified primary estimate uses MAR multiple imputation and Rubin pooling; the complete-outcome regression is a sensitivity analysis. Report raw units, predictions and assignment contrasts at moderator values -1, 0 and +1 SD when supported by data, with their intervals; do not dichotomise PAQ. A spline sensitivity analysis will examine departures from linear moderation. The treatment contrast concerns access to the scaffold regardless of use; no postassignment variables enter the primary model. The observed-outcome fit estimates the full randomised-population target only under appropriate missingness assumptions. Trait scores are not randomised, so no causal claim about changing alexithymia is justified.

\subsection{Secondary and exploratory analyses}
H2 uses post-task clarity regressed on baseline clarity, N-DIF, assignment, initial distress and scenario mode. H3 uses post-task clarity regressed on baseline clarity, assignment, N-DDF and the prespecified baseline covariates, with no postassignment adjustment. Use HC3 and Holm adjustment over the two secondary tests; report both unadjusted estimates and adjusted significance decisions. Exploratory analyses describe agency, completion, suggestion uptake, distress and T2; association of uptake with outcomes is not a causal mediation analysis. Bounded and ordinal outcome models are sensitivity analyses with clearly differentiated estimands. Report model diagnostics, influence and uncertainty without dropping influential points solely to improve significance. The H3 design matrix contains baseline clarity only once. H2 and H3 complete-outcome fits are secondary and require separate missing-outcome sensitivity checks.

\subsection{Missingness and robustness}
Report the participant flow by assigned group, including verified no sketch, withdrawal and item nonresponse. Retain all randomised participants in the denominator. Model availability of the composite using assignment, moderator, their interaction and baseline covariates; report predicted availability quantiles and overlap by arm. Verified structural zeros are retained throughout. For missing composite scores use arm-specific predictive mean matching (five nearest donors, posterior regression draws, 50 imputations), using the primary design predictors. The arm-specific model preserves the assignment-by-moderator interaction and the observed bounded support, including zeros. Pool coefficients and HC3 covariance matrices using Rubin's rules and small-sample degrees of freedom. Report complete-outcome and MAR estimates together. Vary only imputed scores by additive arm-specific and moderator-dependent offsets, clipped to 0--10; use common random numbers across scenarios and report interval and sign changes. The sensitivity grid uses intercept offsets of -2, 0 and 2 points for each arm and differential moderator slopes of -1, 0 and 1 point per SD. These are assumptions, not estimated missing outcomes. An applicable-fit-only model among completers is descriptive because conditioning on postassignment completion can induce selection bias. Report completion contrasts separately; do not infer a direct effect on expression independent of completion. Predefine duplicate rules using study-issued anonymous tokens; no IP or personal platform identifiers are retained.

\section{Ethics and open science}
The task is nonclinical and does not provide therapy, diagnosis or crisis support. Avoid traumatic prompts, permit skipping or stopping without penalty and display institution-approved support information. Recruitment requires ethics review, consent, instrument permission, stimulus rights and institutional data-governance approval. Pre-generated stimuli avoid transmitting personal descriptions to external systems. Any local personal text is discarded at session end. Store only approved pseudonymised numeric records in the institutional environment; publish protocols, approved synthetic examples, analysis code and disclosure-reviewed aggregate results. No human data or access credentials will be placed in the repository. Preregister the final materials, estimands, exclusions, analysis and stopping rule before confirmatory data collection.

\section{Discussion}
The design evaluates a narrow, manipulable component of AI-assisted musical expression. It separates ease of external representation from the user's internal appraisal and permits comparison across the range of alexithymic traits. Its main limitations are an unvalidated task-specific fit rating, self-report bias, constrained assembly instead of unrestricted generation, convenience recruitment and the inability to blind participants. A positive interaction would support differential expressive scaffolding in this task, not treatment efficacy, enduring change or generalisation to all music tools. A future study could use a rights-cleared local generator and independently scored emotion descriptions, with a follow-up beyond 24 hours.

\section{Study status}
This is a prospective protocol. No empirical results, ethics approval, recruitment or registration are asserted. The statistical planning calculation and synthetic Monte Carlo checks are design sensitivity analyses, not results from human participants. The composite endpoint and sensitivity grid must be accepted before preregistration.

\printbibliography
\end{document}
